\chapter{Real Mode}

\section{x86 Real Mode}

This is the legacy 16-bit operating mode that every modern x86 processor starts
up in (backwards compatibility with the original 8086 chip). In this mode, the
CPU has direct, unrestricted access to system memory and hardware, with no
memory protection or virtual memory.

\section{Segmentation}

Because real mode uses 16-bit registers but needs to address more than 64 KB of
memory, it uses a segment-plus-offset scheme. It combines a segment register
(shifted left by 4 bits or simply multiplied by 16) with an offset address to
form a 20-bit physical address, allowing access up to 1MB of memory.

\subsection{The Concept of Segmentation}

Conceptually, segmentation is a way of dividing up memory into smaller,
manageable chunks (called segments) rather treating entire memory space as one
massive, continuous line.
In x86 real mode, the CPU does not just look at memory as a single flat address
space. Instead, it looks through a ``window''. That window is defined by two
pieces working together:

\begin{itemize}
\tightlist
\item
  \textbf{The Segment}: Points to the starting paragraph or block of memory.
\item
  \textbf{The Offset} : Points to the exact by inside that block.
\end{itemize}

The specific way x86 real mode segmentation works (\texttt{Segment\ *\ 16\ +\ Offset}) comes down to a hardware design limitation from the early days of microprocessors:

\begin{enumerate}
\def\labelenumi{\arabic{enumi}.}
\tightlist
\item
  \textbf{The 16-Bit Hardware Constraint:} The original x86 chips were built using 16-bit architecture. This meant the internal registers could naturally hold values up to 64 KB.
\item
  \textbf{The Need for More Memory:} Engineers wanted the computer to access 1 MB of physical memory, which requires a 20-bit address (\(2^{20} = 1,048,576\) bytes).
\item
  \textbf{The Engineering Compromise:} Instead of completely redesigning the chip to use native 20-bit registers (which was impractical and expensive at the time), they found a clever workaround using two 16-bit values.
\end{enumerate}

By taking a 16-bit segment value and shifting it left by 4 bits (which is the same as multiplying by 16, or adding a single hex digit), they turned it into a 20-bit base address. Then, adding a 16-bit offset allowed the CPU to slide that 64KB window anywhere across the 1MB space.

In short, it was a practical hardware hack to squeeze a 20-bit address range out of 16-bit components.

Below table shows \textbf{Implicit Segment Register Usage in x86 Real Mode}.

\begin{longtable}[]{@{}
  >{\raggedright\arraybackslash}p{(\linewidth - 4\tabcolsep) * \real{0.3333}}
  >{\raggedright\arraybackslash}p{(\linewidth - 4\tabcolsep) * \real{0.3333}}
  >{\raggedright\arraybackslash}p{(\linewidth - 4\tabcolsep) * \real{0.3333}}@{}}
\toprule\noalign{}
\begin{minipage}[b]{\linewidth}\raggedright
Operation Type
\end{minipage} & \begin{minipage}[b]{\linewidth}\raggedright
Default Segment Register
\end{minipage} & \begin{minipage}[b]{\linewidth}\raggedright
Description / Typical Usage
\end{minipage} \\
\midrule\noalign{}
\endhead
\bottomrule\noalign{}
\endlastfoot
\textbf{Instruction Fetch} & \texttt{CS} (Code Segment) & Points to the segment containing the executable machine code instructions (\texttt{IP} offset). \\
\textbf{Stack Operations} (\texttt{PUSH}, \texttt{POP}, \texttt{CALL}, \texttt{RET}) & \texttt{SS} (Stack Segment) & Points to the stack area in memory (\texttt{SP} offset). \\
\textbf{General Data Access} (default) & \texttt{DS} (Data Segment) & Used for most general-purpose variable reads and writes (\texttt{BX}, \texttt{SI}, \texttt{DI}, or direct offsets). \\
\textbf{String Operations} (Source) & \texttt{DS} (Data Segment) & Used as the source segment for string manipulation instructions (paired with \texttt{SI}). \\
\textbf{String Operations} (Destination) & \texttt{ES} (Extra Segment) & Strictly required as the destination segment for string instructions like \texttt{MOVS} (paired with \texttt{DI}). \\
\end{longtable}

\section{Memory Operations}

This covers how instructions actually read from and write to that physical
memory using pointers, offsets, and registers (like DS, CS, SS, ES).

\begin{lstlisting}[language=gas]
.intel_syntax noprefix
.code16

_start:
    # 1. set up the stack segment (SS) and stack pointer (SP)
    # whenever PUSH or POP instruction, the CPU automatically
    # uses the stack segment (SS) combined with stack pointer(SP)
    # to find the target memory address.
    # SS:SP --> linear address = SS << 4 + SP
    # for example,
    # mov ax, 0x0000
    # mov ss, ax
    # mov sp, 0x7C00
    # PUSH 0x20
    # the value 0x20 will be in memory address of 0x0000:0x7C00
    # 0x0000 << 4 + 0x7C00 = 0x7C00
    # PUSH 0x20 does SP = SP - 2 (since it's 16-bit real mode)
    # Therefore, SP becomes 0x7BFE
    # the exact logic applies for the POP instruction, just with
    # the inverse stack pointer arithmetic.
    # POP ax --> loads the value from SS:[SP] (or [SS:SP]) into ax
    # and increments the pointer: SP = SP + 2
    mov ax, 0x0000
    mov ss, ax
    mov sp, 0x7c00

    # 2. set up the data segment (DS) for general variables
    mov ax, 0x07c0
    mov ds, ax

    # 3. set up the extra segment (ES)
    mov ax, 0xb800
    mov es, ax

    # 4. code segment (CS)
    # CS is set automatically by the loader/BIOS, but can be targeted
    # with a far jump if you need to change execution segments
    # jmp 0x07c0:load_data

load_data:
    # using DS implicitly for reading a variable
    mov al, [my_variable]

    # using ES explicitly with a segment override to write to video memory
    mov byte [es:0x0000], al
    cli # clear interrupts
    hlt # halt

my_variable db 'A'
\end{lstlisting}

\section{Port I/O and Memory-Mapped I/O}

In x86 real mode, the processor communicates with hardware and peripherals
through two distinct channels.

\begin{itemize}
\tightlist
\item
  \textbf{Memory-Mapped I/O}: Hardware components share the same 1MB RAM space.
  For instance, graphics memory lives right alongside
  regular variables, allowing the processor to interact
  with device simply by using standard data-movement
  instructions like \texttt{MOV}.
\item
  \textbf{Port I/O} : Devices can also reside on an entirely isolated 64KB address
  space. Because regular memory instructions won't work here,
  the CPU relies on specialized commands -- specifically \texttt{IN}
  to pull data from a hardware port into a register, and \texttt{OUT}
  to push data from a register out to a device.
\end{itemize}

\section{Interrupts}

An interrupt is essential an alarm or signal sent to the CPU, telling it to
pause whatever it is currently doing to handle an urgent event or request.

An interrupt service (often called interrupt service routine or ISR) is the
specific block of code designed to handle that event. Once the routine
finishes its job, the CPU picks up right where it left off.

For example, when you press a key on your keyboard, it fires an interrupt
, and the CPU runs a quick keyboard service routine to read the keystroke
before returning to your main program.

\subsection{Interrupts in Real Mode}

In real mode, interrupt handling is simple and direct:

\begin{itemize}
\item
  \textbf{The Interrupt Vector Table (IVT)}: Real mode uses a table located at the very
  bottom of memory (start at address 0x0000). This table contains a list of 256
  pointers (segment and offset addresses) pointing to different interrupt handlers.
\item
  \textbf{How it works}: When an interrupt occurs (or when code run an instruction like
  \texttt{INT\ 0x10}), the CPU looks up the matching number in the IVT, jumps to that
  CS:IP address, runs the routine, and returns.
\item
  \textbf{BIOS/DOS Integration}: Real mode heavily relies on these software interrupts
  built-in BIOS services to read disks, write to the screen, or check the
  keyboard.
\end{itemize}

\subsection{Interrupts in Protected Mode}

Protected mode introduces security, multitasking, and stability, which changes
how interrupts work:

\begin{itemize}
\item
  \textbf{The Interrupt Descriptor Table}: Instead of a simple pointer table like the
  IVT, protected mode uses the IDT.
\item
  \textbf{Descriptor Structures}: Each entry in the IDT is a complex descriptor rather
  than a raw address. It includes a segment selector,
  an offset, access rights, and a descriptor privilege
  level (DPL).
\item
  \textbf{Privilege Checks}: Protected mode ensures that user programs (running in ring
  3) cannot maliciously trigger or hijack system interrupts
  meant only for the kernel (running in ring 0).
\end{itemize}
